\section{Research Question}
\label{sec:doctrine-question}

\subsection{Problem Statement}
\label{sec:doctrine-question-problem}

The foundational problem addressed by the doctrine corpus is this: consequential software
systems presently accumulate activity records downstream of the work that matters. Logs,
dashboards, analytics, and audit trails are collected \emph{after} execution has already
occurred, by systems that had no governing role in that execution. The activity record says
something happened; it does not and cannot prove that what happened was lawful, witnessed,
refusal-capable, or replayable.

The \texttt{BLUE\_RIVER\_DAM.md} doctrine file frames this as the sediment problem:
``downstream data is sediment --- residue left by work. By the time logs, dashboards, and
analytics appear, the governing act already happened.'' Under this regime, governance is
retrospective narration rather than upstream law. Audit trails are constructed post hoc from
records that were never designed to carry evidentiary weight.

This problem is not merely organizational. It is structural. A system that cannot refuse
non-conforming execution cannot govern it. A system whose conformance checks are performed
after execution has no mechanism to stop non-conforming execution from propagating into
downstream claims. Board-level projections, M\&A diligence packages, and regulatory
submissions built on such a foundation carry claims that cannot be independently replayed or
falsified. The doctrine corpus addresses this structural defect at its root by specifying
what an upstream admission and refusal layer must look like in formal, type-theoretic terms.

\subsection{Old-Regime Assumption Challenged}
\label{sec:doctrine-question-old-regime}

The old regime assumes the following sequence to be sufficient and correct: write code,
deploy, observe execution, collect logs, and explain observed behavior in retrospect. This
assumption underlies conventional business intelligence, observability tooling, workflow
automation, and monitoring dashboards. Each of these disciplines records what happened but
does not adjudicate whether what happened conforms to what was declared.

\texttt{PROCESS\_INTELLIGENCE\_DEFINED.md} catalogues the specific distinctions:

\begin{itemize}
  \item Process intelligence is \emph{not} business intelligence. BI aggregates reported
    data. Process intelligence mines actual execution behavior and detects model-log
    divergence.
  \item Process intelligence is \emph{not} observability. Observability records what
    happened. Process intelligence proves whether what happened conforms to what was declared.
  \item Process intelligence is \emph{not} workflow automation. Automation executes declared
    steps. Process intelligence validates that execution conforms to law and refuses
    non-conforming paths.
  \item Process intelligence is \emph{not} monitoring dashboards. Dashboards display metrics.
    Process intelligence issues typed verdicts grounded in formal process models.
\end{itemize}

The challenged assumption is that recording activity is equivalent to governing it. The
doctrine corpus argues this assumption produces unclosed process truth: activity records with
no admission gate, no refusal surface, no witness, and no replayable receipt. Claims built
on such records are inadmissible under the board-admissibility criteria established in the
M\&A doctrine.

\subsection{Hypothesis}
\label{sec:doctrine-question-hypothesis}

The hypothesis of the doctrine corpus is:

\begin{quote}
Full-lifecycle process intelligence --- defined as the systematic application of process
mining, conformance checking, and evidence-grounded execution control from design through
decommission --- can be manufactured from a typed, receipt-bearing, refusal-capable upstream
admission layer such that every downstream claim is independently replayable and every gap is
explicitly accounted for rather than silently absorbed.
\end{quote}

This hypothesis has several formal components. First, the admission layer must be
type-enforced: \texttt{Evidence\textless{}T, State, W\textgreater{}} with
\texttt{PhantomData} lifecycle states and witness tags ensures that only lawfully admitted
evidence can enter execution paths, and this guarantee is enforced at compile time, not at
runtime. Second, refusals must be named: generic string-typed catch-alls such as
\texttt{ValidationError(String)} are explicitly forbidden by doctrine, and every refusal must
name the specific violated law. Third, every critical transition must emit a
\texttt{Receipt\textless{}T, W\textgreater{}} --- a typed, witnessed proof artifact --- so
that the audit chain is machine-verifiable rather than narrative. Fourth, gaps must be
documented rather than hidden: the hostile-witness posture of the research program, evidenced
by the self-discovered \texttt{GAP\_001} finding, demonstrates that honest accounting of
structural defects is an integral part of the doctrine, not a failure mode.

\subsection{Success Criteria}
\label{sec:doctrine-question-criteria}

The doctrine corpus defines success using the ALIVE gate framework established in
\texttt{PROCESS\_INTELLIGENCE\_ALIVE\_001.md}. The sealed checkpoint at 588 commits with
12/12 criteria verified constitutes the primary success evidence. The six specific gate
criteria relevant to the doctrine layer are:

\begin{enumerate}
  \item \textbf{Admissibility boundary}: $R \vdash P_i = \mu(O^*, T, L)$ --- compile-time
    state-space bounds are validated.
  \item \textbf{Autonomic actuation}: $\alpha(K, P, L, T) \to \tau$ --- typestate
    compile-fail loops preventing invalid state transitions are verified.
  \item \textbf{Token-game fitness}: $f(L, N) = 1 - \frac{\sum m}{\sum c} - \frac{\sum
    r}{\sum p}$ --- verified via automated execution gate.
  \item \textbf{OCPQ refinement}: $p \sqsubseteq_L c$ --- subset inclusion and matching
    bindings are programmatically verified.
  \item \textbf{Decommissioning retirement}: $\delta(P)$ --- formal closure maps preventing
    orphan dependencies and active leaks.
  \item \textbf{Downstream authorization}: six specific workflows authorized and receipted
    against the sealed checkpoint.
\end{enumerate}

The doctrine corpus also defines partial success honestly. The overall project status is
\textbf{PARTIAL}: the doctrine layer is sealed and authoritative, but the downstream
execution layer (wasm4pm) has not yet implemented the bridge to receive admitted evidence
from wasm4pm-compat. All fifteen algorithm families in the taxonomy show wasm4pm execution as
\emph{missing}. PARTIAL is declared explicitly, not absorbed silently. This is itself a
success criterion: honest accounting of incompleteness is a core requirement of the doctrine.
